\section{Closed geometry}
\label{sec:octahedron}
The three complex directions and eight connection components from Section~\ref{sec:color} now give the axes and face labels of an octahedron. Closing this representation supplies the patch arrangement; placing the actual source on those patches then gives the smooth field and its separating boundary.

\subsection{The closed representation}
Let $e_1,e_2,e_3$ be the standard basis of the three-dimensional drawing space and choose
\CurvatureSignedEndpointsDisplay
with drawing scale $\DisplaySize>0$. Taking one signed endpoint from each axis produces
\CurvatureOctahedralFacesDisplay
The eight sign triples give eight faces. Every edge belongs to two faces with opposite induced orientations, so their boundary contributions cancel,
\CurvatureClosedSurfaceDisplay

This closed arrangement also fixes the total angular defect. At a vertex $v$, let $\delta_v$ be a full turn minus the incident face angles. Their sum is
\CurvatureTotalDefectDisplay
Local defects can redistribute under nondegenerate piecewise Euclidean deformations that preserve the connectivity, but their total remains fixed. Appendix~\ref{app:source-reduction} supplies the count. This angular result belongs to the representation; the smooth field's curvature is obtained from its metric after assembly.

\begin{figure}[!htbp]
\centering
\resizebox{.62\linewidth}{!}{\CarrierScene{3}{3}{2}}
\FaceKey{2}
\caption{Closed representation of the eight color components, with six off-diagonal and two diagonal face labels. Shared edges cancel, and the total angular defect is independent of the drawing scale $\DisplaySize$.}
\label{fig:complete}
\end{figure}
\FloatBarrier

\subsection{Assembling the smooth field}
Radial projection turns each face into a patch on the unit sphere. For a sign triple $s=(s_1,s_2,s_3)$, let $\lambda_i\ge0$ locate a point in the face, with $\lambda_1+\lambda_2+\lambda_3=1$. Its unit direction is
\ClosedPatchDirectionDisplay
A vanishing weight places the point on a shared edge, where the adjacent descriptions agree. At a vertex a smooth chart contains the four incident sectors. Together these patches cover the sphere.

To put the constructed field on these patches, use its original four source vectors $b_0,b_1,b_2,b_3$. The columns of $\mathsf D(q)$ record $b_1-b_0,b_2-b_0,b_3-b_0$ in source coordinates $q=(x,y,z)$. Their weighted differences undergo the vector coupling and assemble the self-adjoint matrices $X_p$; the same matrix coupling then gives $Y=\Lambda(X)$ and $Z=\Lambda(Y)$. Thus $\mathsf D$ describes the input to the response chain of Sections~\ref{sec:coupling}--\ref{sec:color}. Appendix~\ref{app:source-coordinates} specifies the weights and assembly.

With reference point $q_*=(-1,-1,1)$, the point $q_*+r\widehat r_s(\lambda)$ selects this source on a sphere of coordinate radius $r>0$. Each patch carries the resulting $X,Y,Z$ and their induced geometry.

These matrices need a common source scale. Let $R(\mathsf D)$ be the positive scale function determined by the coefficient geometry and set $R_0=R(\mathsf D(q_*))$. The normalized patch is
\ClosedSourcePatchDisplay
The exponent follows from $R(a\mathsf D)=a^{324}R(\mathsf D)$ for $a>0$. This factor changes the scale of the matrix while leaving the selected source-coordinate point fixed. Radius $r$ locates the surface; $R_0$ specifies the common scale of its data.

Each patch is a restriction of the same smooth source map. Source values and derivatives therefore agree at shared edges and vertices after the corresponding coordinate changes. Applying the coupling gives a smooth field, metric and connection across the assembled surface.

\subsection{The boundary and its two sides}
On the fixed-scale spatial slice $\Sigma_R$, the image of the coordinate ball $B_r(q_*)$ defines $\Sigma_-$. Its boundary is $S_r=\partial\Sigma_-$, and the complementary region is $\Sigma_+=\Sigma_R\setminus\overline{\Sigma_-}$. The signs label sides of the surface. The radius remains general within the regular neighborhood of the source.

Allowing the source coordinates and scale to vary continues this geometry across $S_r$. Neighborhoods on its two sides, called collars, have matching induced metrics and second fundamental forms when compared with a common normal. Thus the separating surface divides a matched geometry with source data available on both sides.

\FloatBarrier
